\documentclass[a4paper,11pt]{amsart}

\usepackage[margin=1in]{geometry}
\usepackage[T1]{fontenc}
\usepackage{amsmath,amssymb,mathtools}
\usepackage[protrusion=true,expansion=false]{microtype}

\newtheorem{theorem}{Theorem}[section]
\newtheorem{proposition}[theorem]{Proposition}
\newtheorem{lemma}[theorem]{Lemma}
\newtheorem{definition}[theorem]{Definition}
\theoremstyle{remark}
\newtheorem{remark}[theorem]{Remark}

\newcommand{\Z}{\mathbb Z}
\newcommand{\Q}{\mathbb Q}
\newcommand{\R}{\mathbb R}
\newcommand{\C}{\mathbb C}
\newcommand{\PP}{\mathbb P}
\DeclareMathOperator{\Comp}{Comp}
\DeclareMathOperator{\GL}{GL}
\DeclareMathOperator{\Ogrp}{O}
\DeclareMathOperator{\SO}{SO}
\DeclareMathOperator{\Lie}{Lie}
\DeclareMathOperator{\rank}{rank}
\DeclareMathOperator{\im}{im}

\numberwithin{equation}{section}

\title{Odd All-Half Thinness}
\date{}

\begin{document}

\begin{abstract}
We prove thinness for the odd all-half hypergeometric tower
\[
\Gamma_d=\left\langle \Comp((x-1)^d),\Comp((x+1)^d)\right\rangle,
\qquad d\ge5\text{ odd}.
\]
After a rational Levelt normal-form change, the group is generated by a regular unipotent isometry and a reflection.  The main input is a positivity theorem for an upper-Hessenberg determinant attached to the even kernel coming from \(((1+z)/(1-z))^d\).  This positivity gives a projective free-product certificate
\[
\PP\Gamma_d\cong (\Z/2\Z)*\Z.
\]
Zariski density follows from the conjugates of the reflection, and infinite index follows from an elementary abelian-subgroup obstruction.
\end{abstract}

\maketitle

\section{The normal form}

Let \(d=2M+1\ge5\) be odd and set
\[
        N=d-1=2M.
\]
Define \(c_d(n)\), \(n\in\Z\), by
\[
        \sum_{n\ge0}c_d(n)z^n
        =
        \left(\frac{1+z}{1-z}\right)^d,
        \qquad
        c_d(n)=0\quad(n<0).
\]
Define the even kernel \(b_d:\Z\to\Z\) by
\[
        b_d(0)=2,
        \qquad
        b_d(n)=c_d(|n|)\quad(n\ne0).
\]
Throughout, \(\Comp(h)\) denotes the column companion matrix: if
\[
        h(x)=x^d+a_{d-1}x^{d-1}+\cdots+a_1x+a_0,
\]
then \(\Comp(h)e_i=e_{i+1}\) for \(0\le i<d-1\), and
\[
        \Comp(h)e_{d-1}=-\sum_{i=0}^{d-1}a_i e_i.
\]
For \(m\in\Z\), put
\[
        \rho_m(t)=\left(\frac{1+t}{1-t}\right)^m
\]
as a formal power series at \(t=0\).  Also define
\[
        \Lambda_d(F)=2\sum_{j=0}^{M}[t^{2j}]F(t).
\]
When \(d\) is fixed, we write simply \(\Lambda=\Lambda_d\).

Let \(V\) be a \(d\)-dimensional rational vector space with basis
\[
        e_0,e_1,\ldots,e_{d-1}.
\]
Define \(A\in\GL(V)\) by the companion rule for \((x-1)^d\):
\[
        Ae_i=e_{i+1}\quad(0\le i<d-1),
\]
and
\[
        Ae_{d-1}
        =
        \sum_{j=0}^{d-1}(-1)^j\binom dj e_j.
\]
Thus \((A-I)^d=0\), and \(A\) is one regular unipotent Jordan block.

Define a symmetric bilinear form \(Q=Q_d\) on \(V\) by
\[
        Q(e_i,e_j)=b_d(j-i)
        \qquad(0\le i,j\le d-1).
\]
Since \(b_d\) is even, this is symmetric.  This is a normal-form construction: at this point we are not claiming that the standard companion basis carries this Toeplitz Gram matrix.  The passage back to companion coordinates is made only after Proposition~\ref{prop:normal-form}, using the rank-one Levelt uniqueness lemma and a commensurability argument for the resulting lattices.

Finally define
\[
        R(x)=x-Q(x,e_0)e_0.
\]
Since \(Q(e_0,e_0)=b_d(0)=2\), this map will be the orthogonal reflection in \(e_0^\perp\) once nondegeneracy of \(Q\) has been proved.

\begin{lemma}\label{lem:b-polynomial}
The function \(b_d(n)\) is the restriction to \(\Z\) of an even polynomial \(B_d(n)\in\Q[n]\) of degree \(d-1\), with leading coefficient
\[
        \frac{2^d}{(d-1)!}.
\]
Consequently,
\[
        \sum_{r=0}^{d}(-1)^r\binom dr b_d(n-r)=0
        \qquad(n\in\Z).
\]
\end{lemma}

\begin{proof}
With the notation just introduced, a residue computation after the change of variables
\[
        z=\frac{1-t}{1+t}
\]
gives
\begin{equation}\label{eq:b-Lambda-rho}
        b_d(n)=\Lambda(\rho_n)
        \qquad(n\in\Z).
\end{equation}
Indeed, for \(n>0\),
\[
\begin{aligned}
        c_d(n)
        &=-\operatorname*{Res}_{z=1}
        \left(\frac{1+z}{1-z}\right)^d z^{-n-1}\,dz        \\
        &=2[t^{d-1}]\frac{\rho_n(t)}{1-t^2}
        =2\sum_{j=0}^{M}[t^{2j}]\rho_n(t).
\end{aligned}
\]
For \(n=0\), \(\Lambda(1)=2=b_d(0)\).  For negative \(n\), use \(\rho_{-n}(t)=\rho_n(-t)\) and the fact that \(\Lambda\) reads only even coefficients.

For fixed \(0\le r\le N\), the coefficient \([t^r]\rho_n(t)\) is a polynomial in \(n\) of degree at most \(r\), because
\[
        \rho_n(t)=\exp\left(n\log\frac{1+t}{1-t}\right).
\]
Since \(\Lambda\) uses only coefficients up to degree \(N=d-1\), \(b_d(n)\) is a polynomial in \(n\) of degree at most \(d-1\).  The equality \(b_d(-n)=b_d(n)\) makes this polynomial even.  Its leading term comes only from the coefficient of \(t^N\), and
\[
        \log\frac{1+t}{1-t}=2t+O(t^3),
\]
so
\[
        [t^N]\rho_n(t)=\frac{2^N}{N!}n^N+O(n^{N-1}).
\]
Multiplication by the outer factor \(2\) in \(\Lambda\) gives leading coefficient \(2^{N+1}/N!=2^d/(d-1)!\).  The finite-difference identity is the vanishing of the \(d\)-th finite difference of a polynomial of degree \(d-1\).
\end{proof}

\begin{lemma}[Levelt uniqueness in the rank-one case]\label{lem:levelt-uniqueness}
Let \(f,g\in\Q[x]\) be coprime monic polynomials of degree \(d\).  Let \((A,B)\) and \((A',B')\) be pairs of rational linear maps on \(d\)-dimensional rational vector spaces such that
\[
        \chi_A=\chi_{A'}=f,
        \qquad
        \chi_B=\chi_{B'}=g,
        \qquad
        \rank(B-A)=\rank(B'-A')=1.
\]
Then the two pairs are rationally conjugate.
\end{lemma}

\begin{proof}
It is enough to put a single pair \((A,B)\) into a basis depending only on \(f\) and \(g\).  Choose
\[
        0\ne u\in\im(B-A),
\]
and write
\[
        B=A+u\otimes\varphi,
        \qquad
        (u\otimes\varphi)(x)=\varphi(x)u
\]
for a rational covector \(\varphi\).  Let
\[
        W=\operatorname{span}_{\Q}\{A^i u:i\ge0\}.
\]
Then \(W\) is \(A\)-invariant, and it is also \(B\)-invariant because \((B-A)V\subseteq \Q u\subseteq W\).  On \(V/W\), the induced maps of \(A\) and \(B\) are equal.  Hence the characteristic polynomial of this common quotient map divides both \(f\) and \(g\).  Since \((f,g)=1\), the quotient is zero.  Thus
\[
        u,Au,\ldots,A^{d-1}u
\]
is a basis of \(V\).

In this cyclic basis, the matrix of \(A\) is the column companion matrix of \(f\).  The matrix of \(B\) is determined by the scalars
\[
        \alpha_i=\varphi(A^i u),
        \qquad 0\le i\le d-1.
\]
The characteristic polynomial of \(B\) determines these scalars.  Indeed, the matrix determinant lemma gives, as a Laurent expansion at \(\lambda=\infty\),
\[
\begin{aligned}
        \frac{g(\lambda)}{f(\lambda)}
        &=
        \frac{\det(\lambda I-B)}{\det(\lambda I-A)}      \\
        &=1-\varphi\bigl((\lambda I-A)^{-1}u\bigr)       \\
        &=1-\sum_{i=0}^{d-1}\alpha_i\lambda^{-i-1}
          +O(\lambda^{-d-1}).
\end{aligned}
\]
Thus \(\alpha_0,\ldots,\alpha_{d-1}\) are the first \(d\) coefficients of the Laurent series of \(1-g/f\) at infinity.  They depend only on \(f\) and \(g\).  Therefore every pair satisfying the hypotheses has the same matrices in its cyclic basis, and any two such pairs are rationally conjugate.
\end{proof}

\begin{proposition}\label{prop:orthogonal-normal-form}
The form \(Q\) is nondegenerate, and both \(A\) and \(R\) preserve \(Q\).
\end{proposition}

\begin{proof}
We first prove that \(A\) preserves \(Q\).  If \(i,j<N\), then
\[
        Q(Ae_i,Ae_j)=Q(e_{i+1},e_{j+1})=b_d(j-i)=Q(e_i,e_j).
\]
If \(j<N\), then Lemma~\ref{lem:b-polynomial} gives
\[
\begin{aligned}
        Q(Ae_N,Ae_j)
        &=
        \sum_{r=0}^{N}(-1)^r\binom dr\, b_d(j+1-r)       \\
        &=
        b_d(j-N)
        =
        Q(e_N,e_j).
\end{aligned}
\]
The equality with \(e_j,e_N\) follows by symmetry.  Finally set
\[
        S=
        \sum_{r,s=0}^{N}
        (-1)^{r+s}\binom dr\binom ds\,b_d(s-r).
\]
The full double sum with \(0\le r,s\le d\) is zero, by applying the finite-difference identity in the \(s\)-variable and using the evenness of \(b_d\).  Also
\[
        \sum_{r=0}^{N}(-1)^r\binom dr b_d(d-r)=b_d(0),
\]
because the missing \(r=d\) term in the corresponding full finite difference is \((-1)^d b_d(0)=-b_d(0)\).  Therefore the two boundary sums with one index equal to \(d\) each contribute \(-b_d(0)\), while the corner \((r,s)=(d,d)\) contributes \(b_d(0)\).  Hence the full double sum is \(S-b_d(0)\), so \(S=b_d(0)\).  This proves
\[
        Q(Ae_N,Ae_N)=Q(e_N,e_N).
\]
Thus \(A^tQA=Q\).

The formula for \(R\) and the equality \(Q(e_0,e_0)=2\) give, for all \(x,y\in V\),
\[
\begin{aligned}
        Q(Rx,Ry)
        &=Q(x,y)-Q(x,e_0)Q(e_0,y)-Q(y,e_0)Q(x,e_0)  \\
        &\qquad +Q(x,e_0)Q(y,e_0)Q(e_0,e_0)          \\
        &=Q(x,y).
\end{aligned}
\]
So \(R\) also preserves \(Q\).

It remains to prove nondegeneracy.  Let \(\mathcal R\) be the radical of \(Q\).  Since \(A\) preserves \(Q\), the radical is \(A\)-stable.  If \(\mathcal R\ne0\), then, because \(A\) is a single unipotent Jordan block, \(\mathcal R\) contains the fixed line of \(A\), namely the line spanned by
\[
        \eta=(A-I)^N e_0.
\]
Indeed, if \(0\ne w\in\mathcal R\), choose \(m\) maximal with \((A-I)^mw\ne0\); then \((A-I)^mw\) is fixed by \(A\), and the fixed space of a single Jordan block is one-dimensional.
But
\[
\begin{aligned}
        Q(e_0,\eta)
        &=
        \sum_{r=0}^{N}(-1)^{N-r}\binom Nr b_d(r)      \\
        &=
        N!\,[n^N]B_d(n)
        =2^d\ne0,
\end{aligned}
\]
where \(B_d\) is the polynomial from Lemma~\ref{lem:b-polynomial}.  Thus \(\eta\) is not in the radical, a contradiction.  Therefore \(Q\) is nondegenerate.
\end{proof}

\begin{lemma}\label{lem:orbit-gram}
For all \(p,q\in\Z\),
\[
        Q(A^p e_0,A^q e_0)=b_d(q-p).
\]
\end{lemma}

\begin{proof}
By \(A\)-invariance of \(Q\), it is enough to prove
\[
        Q(e_0,A^n e_0)=b_d(n)
        \qquad(n\in\Z).
\]
For \(n\ge0\), put \(s(n)=Q(e_0,A^n e_0)\).  The equality \((A-I)^d=0\) implies that \(s(n)\) satisfies the forward finite-difference equation
\[
        \sum_{r=0}^{d}(-1)^{d-r}\binom dr s(n+r)=0.
\]
Since \(b_d\) is the restriction of a polynomial of degree \(d-1\), it satisfies the same forward equation.  The two sequences agree for \(0\le n\le d-1\) by the definition of \(Q\).  Hence \(s(n)=b_d(n)\) for all \(n\ge0\).  For \(n<0\), write \(n=-m\) with \(m>0\).  Then
\[
        Q(e_0,A^{-m}e_0)=Q(A^m e_0,e_0)=b_d(m)=b_d(-m),
\]
using \(A\)-invariance, symmetry of \(Q\), and evenness of \(b_d\).
\end{proof}

\begin{proposition}\label{prop:normal-form}
Let \(B=RA\).  Then
\[
        \chi_A(x)=(x-1)^d,
        \qquad
        \chi_B(x)=(x+1)^d,
        \qquad
        (B+I)^d=0.
\]
Moreover, the pair \((A,B)\) is rationally conjugate to the companion pair
\[
        \left(\Comp((x-1)^d),\Comp((x+1)^d)\right).
\]
\end{proposition}

\begin{proof}
The assertion for \(A\) is part of its definition.  For \(B=RA\), observe that
\[
        B=A-e_0\otimes\bigl(x\mapsto Q(Ax,e_0)\bigr),
\]
so \(B\) is a rank-one perturbation of \(A\).  The matrix determinant lemma gives
\[
\begin{aligned}
        \det(\lambda I-B)
        &=
        \det(\lambda I-A)
        \left(
        1+Q(A(\lambda I-A)^{-1}e_0,e_0)
        \right).
\end{aligned}
\]
As a Laurent expansion at infinity,
\[
\begin{aligned}
        Q(A(\lambda I-A)^{-1}e_0,e_0)
        &=
        \sum_{m\ge1} Q(A^m e_0,e_0)\lambda^{-m}       \\
        &=
        \sum_{m\ge1} b_d(m)\lambda^{-m}
        =
        \sum_{m\ge1} c_d(m)\lambda^{-m},
\end{aligned}
\]
by Lemma~\ref{lem:orbit-gram}.  Therefore
\[
\begin{aligned}
        1+Q(A(\lambda I-A)^{-1}e_0,e_0)
        &=
        \sum_{m\ge0}c_d(m)\lambda^{-m}                  \\
        &=
        \left(\frac{1+\lambda^{-1}}{1-\lambda^{-1}}\right)^d
        =
        \left(\frac{\lambda+1}{\lambda-1}\right)^d.
\end{aligned}
\]
Since both sides are rational functions of \(\lambda\), equality of their Laurent expansions at infinity gives equality as rational functions.  Since \(\det(\lambda I-A)=(\lambda-1)^d\), we obtain
\[
        \det(\lambda I-B)=(\lambda+1)^d.
\]
Cayley-Hamilton then gives \((B+I)^d=0\).

The difference \(B-A=(R-I)A\) has rank one.  The two column companion matrices of \((x-1)^d\) and \((x+1)^d\) also differ by a rank-one matrix, since they differ only in the last companion column.  The polynomials \((x-1)^d\) and \((x+1)^d\) are coprime.  Lemma~\ref{lem:levelt-uniqueness} therefore applies and gives the claimed rational conjugacy.
\end{proof}

\begin{remark}
All subsequent arguments take place in the normal form \(\langle A,R\rangle\).  The final transfer to the companion lattice is by the rational conjugacy in Proposition~\ref{prop:normal-form}.  Clearing denominators in the transported form gives an integral orthogonal group.
\end{remark}

\section{The determinant theorem}

Let
\[
        p_0=0,
        \qquad
        p_j-p_{j-1}=k_j\ne0
        \quad(1\le j\le\ell)
\]
be a finite walk on \(\Z\).  Define the \(\ell\times\ell\) upper-Hessenberg matrix
\[
        H_b(k)=\bigl(h_{ij}\bigr)_{1\le i,j\le \ell}
\]
by
\[
        h_{ij}
        =
        \begin{cases}
        b_d(p_j-p_{i-1}),&i\le j,\\
        1,&i=j+1,\\
        0,&i\ge j+2.
        \end{cases}
\]
Again, \(b_d(p_j-p_{i-1})=b_d(|p_j-p_{i-1}|)\).

\begin{theorem}[Odd walk positivity]\label{thm:walk-positive}
For the fixed odd integer \(d\ge5\), every \(\ell\ge1\), and every finite walk with nonzero steps,
\[
        \det H_b(k)>2.
\]
In particular, \(\det H_b(k)>0\).
\end{theorem}

\begin{proof}
For a formal power series \(F(t)\), write \([F]_{\le N}\) for its truncation to degrees \(0,1,\ldots,N\).  Let
\[
        D_s=\det H_b(k_1,\ldots,k_s),
        \qquad
        D_0=1.
\]
The upper-Hessenberg expansion gives
\begin{equation}\label{eq:hessenberg-rec}
        D_s
        =
        \sum_{i=0}^{s-1}(-1)^{s-1-i}D_i\,b_d(p_s-p_i).
\end{equation}
Define
\[
        P_s(z)=\sum_{i=0}^{s}(-1)^{s-i}D_i z^{p_i},
\]
and
\[
        Q_s(t)=\left[z(t)^{-p_s}P_s(z(t))\right]_{\le N},
        \qquad
        z(t)=\frac{1-t}{1+t}.
\]
Since \(z(t)^{-k}=\rho_k(t)\), we have
\[
        z(t)^{-p_s}P_{s-1}(z(t))
        =
        \rho_{k_s}(t)\,z(t)^{-p_{s-1}}P_{s-1}(z(t)).
\]
Here \(P_s\) is in general a Laurent polynomial, but \(z(t)\) has constant term \(1\), so every integral power \(z(t)^m\) is a formal power series with no negative powers of \(t\).  Also every \(\rho_k(t)=z(t)^{-k}\) has no negative powers.  Since \(\Lambda\) depends only on coefficients of degrees at most \(N\), coefficients of degree at most \(N\) in
\[
        \rho_{k_s}(t) z(t)^{-p_{s-1}}P_{s-1}(z(t))
\]
depend only on the truncation \(Q_{s-1}\).  Therefore the determinant recurrence becomes
\[
        D_s=\Lambda(\rho_{k_s}Q_{s-1}).
\]
Moreover
\[
        P_s(z)=D_s z^{p_s}-P_{s-1}(z),
\]
so, after multiplying by \(z(t)^{-p_s}\) and taking coefficients through degree \(N\),
\begin{equation}\label{eq:Q-recursion}
        D_s=\Lambda(\rho_{k_s}Q_{s-1}),
        \qquad
        Q_s=[D_s-\rho_{k_s}Q_{s-1}]_{\le N}.
\end{equation}

We now prove that this recursion preserves a finite coefficient cone.  Write
\[
        Q(t)=q_0+q_1t+\cdots+q_Nt^N.
\]
For \(\sigma\in\{\pm1\}\), say that \(Q\) is \(\sigma\)-admissible if
\[
        u_0=q_0,
        \qquad
        u_n=-\sigma^nq_n\quad(1\le n\le N)
\]
satisfy
\begin{equation}\label{eq:admissible}
        u_n\ge0\quad(1\le n\le N),
        \qquad
        u_0>\sum_{n=1}^{N}u_n.
\end{equation}

Let \(Q\) be \(\sigma\)-admissible, and let
\[
        k=\tau a,
        \qquad
        \tau\in\{\pm1\},
        \qquad
        a\ge1.
\]
Write
\[
        \rho_a(t)=\sum_{n\ge0}C_a(n)t^n.
\]
The coefficients \(C_a(n)\) are positive from the product expansion \(\rho_a=(1+t)^a(1-t)^{-a}\).  They are nondecreasing because
\[
        (1-t)\rho_a(t)=\frac{(1+t)^a}{(1-t)^{a-1}}
\]
has nonnegative coefficients.  Also the alternating sums
\[
        T_a(m)=(-1)^m\sum_{n=0}^{m}(-1)^nC_a(n)
\]
are positive and nondecreasing.  Indeed,
\[
        \sum_{m\ge0}T_a(m)t^m
        =
        \frac{\rho_a(t)}{1+t}
        =
        \frac{(1+t)^{a-1}}{(1-t)^a},
\]
which has positive coefficients, and multiplication by \(1-t\) gives
\[
        \frac{(1+t)^{a-1}}{(1-t)^{a-1}},
\]
which has nonnegative coefficients.  The coefficient of \(t^m\) is \(T_a(m)-T_a(m-1)\), with \(T_a(-1)=0\), so the sequence \(T_a(m)\) is nondecreasing.

Set
\[
        A_Q=[\rho_kQ]_{\le N}=\sum_{n=0}^{N}a_nt^n,
        \qquad
        V_n=\tau^n a_n.
\]
Using \(q_0=u_0\) and \(q_j=-\sigma^ju_j\) for \(j\ge1\), we get
\begin{equation}\label{eq:Vn}
        V_n
        =
        C_a(n)u_0
        -
        \sum_{j=1}^{n}(\tau\sigma)^jC_a(n-j)u_j.
\end{equation}
Since \(C_a(n-j)\le C_a(n)\),
\begin{equation}\label{eq:Vnpositive}
        V_n
        \ge
        C_a(n)\left(u_0-\sum_{j=1}^{n}u_j\right)>0.
\end{equation}
Thus
\[
        D:=\Lambda(A_Q)=2\sum_{j=0}^{M}a_{2j}
        =
        2\sum_{j=0}^{M}V_{2j}
        >0.
\]

Now define
\[
        Q'=[D-A_Q]_{\le N}.
\]
For \(n\ge1\),
\[
        -\tau^nq'_n=\tau^na_n=V_n>0.
\]
For the dominance condition, note that
\[
\begin{aligned}
        q'_0-\sum_{n=1}^{N}(-\tau^nq'_n)
        &=
        (D-a_0)-\sum_{n=1}^{N}V_n                 \\
        &=
        \sum_{n=0}^{N}(-1)^nV_n.
\end{aligned}
\]
Let \(\varepsilon=\tau\sigma\).  Since \(N=d-1\) is even, \eqref{eq:Vn} gives
\[
        \sum_{n=0}^{N}(-1)^nV_n
        =
        u_0T_a(N)
        -
        \sum_{j=1}^{N}\varepsilon^ju_jT_a(N-j).
\]
Because \(0<T_a(N-j)\le T_a(N)\),
\[
        \sum_{n=0}^{N}(-1)^nV_n
        \ge
        T_a(N)\left(u_0-\sum_{j=1}^{N}u_j\right)>0.
\]
Hence \(Q'\) is \(\tau\)-admissible.

We start with \(Q_0=1\), which is both \(+\)- and \(-\)-admissible.  Applying the preceding paragraph inductively to \eqref{eq:Q-recursion} proves \(D_s>0\) for all \(s\).  Thus \(\det H_b(k)=D_\ell>0\).

Since \(d\ge5\), we have \(N\ge4\).  This strict strengthening is needed later to distinguish a nontrivial reduced word from the scalar \((-1)^\ell I\).  The recursion preserves integral coefficients: \(Q_0=1\in\Z[t]\), and if \(Q\in\Z[t]\), then \([\rho_k Q]_{\le N}\in\Z[t]\).  Indeed, for \(a\ge1\),
\[
        \rho_a(t)=\frac{(1+t)^a}{(1-t)^a}\in\Z[[t]],
        \qquad
        \rho_{-a}(t)=\frac{(1-t)^a}{(1+t)^a}\in\Z[[t]],
\]
so \(\rho_k\in\Z[[t]]\) for every \(k\in\Z\).  Hence \(D=\Lambda([\rho_kQ]_{\le N})\in\Z\), and \(Q'=[D-\rho_kQ]_{\le N}\in\Z[t]\).  Thus all \(V_n\), in particular \(V_0\) and \(V_2\), are positive integers.  Therefore at every step \(s\ge1\),
\[
        D_s=2\sum_{j=0}^{M}V_{2j}\ge2(V_0+V_2)\ge4.
\]
Thus \(D_s>2\).
\end{proof}

\section{Projective freeness}

Let
\[
        G_d=\langle A,R\rangle.
\]
The projective order of \(R\) is \(2\), and the projective order of \(A\) is infinite.  Therefore there is a natural surjective homomorphism
\[
        (\Z/2\Z)*\Z\longrightarrow \PP G_d.
\]

A cyclically reduced mixed word has the form
\[
        w=r a^{k_1} r a^{k_2}\cdots r a^{k_\ell},
        \qquad
        k_i\in\Z\setminus\{0\}.
\]
Let
\[
        p_0=0,
        \qquad
        p_j=k_1+\cdots+k_j.
\]
Let \(W\) be the corresponding matrix:
\[
        W=R A^{k_1}R A^{k_2}\cdots R A^{k_\ell}.
\]

\begin{lemma}[Reflection determinant]\label{lem:reflection-det}
With the notation above,
\[
        Q(e_0,We_0)=(-1)^\ell\det H_b(k_1,\ldots,k_\ell).
\]
\end{lemma}

\begin{proof}
Put \(v_i=A^{p_i}e_0\).  The conjugate
\[
        A^{p}RA^{-p}
\]
is the reflection in the vector \(A^p e_0\), because \(A\) preserves \(Q\).  Hence
\[
        W=(A^{p_0}RA^{-p_0})(A^{p_1}RA^{-p_1})\cdots(A^{p_{\ell-1}}RA^{-p_{\ell-1}})A^{p_\ell}.
\]
Let \(\rho_i=A^{p_i}RA^{-p_i}\), so
\[
        \rho_i(x)=x-Q(x,v_i)v_i.
\]
Since \(\rho_0\) sends \(v_0=e_0\) to \(-v_0\),
\[
        Q(e_0,We_0)
        =Q(v_0,\rho_0\rho_1\cdots\rho_{\ell-1}v_\ell)
        =-Q(v_0,\rho_1\cdots\rho_{\ell-1}v_\ell).
\]
Now expand \(\rho_m=I-P_m\), \(1\le m\le\ell-1\), where
\[
        P_m(x)=Q(x,v_m)v_m.
\]
Choosing the projections with indices
\[
        0<i_1<\cdots<i_{s-1}<\ell
\]
contributes
\[
        -(-1)^{s-1}
        Q(v_0,v_{i_1})Q(v_{i_1},v_{i_2})\cdots Q(v_{i_{s-1}},v_\ell),
\]
with the evident interpretation when no interior index is chosen.  Therefore
\[
        Q(e_0,We_0)
        =
        \sum_{0=i_0<i_1<\cdots<i_s=\ell}
        (-1)^s
        \prod_{a=0}^{s-1}Q(v_{i_a},v_{i_{a+1}}).
\]
By Lemma~\ref{lem:orbit-gram}, this becomes
\[
        Q(e_0,We_0)
        =
        \sum_{0=i_0<i_1<\cdots<i_s=\ell}
        (-1)^s
        \prod_{a=0}^{s-1}b_d(p_{i_{a+1}}-p_{i_a}).
\]
The same interval-chain expansion for the upper-Hessenberg determinant gives
\[
        \det H_b(k)
        =
        \sum_{0=i_0<i_1<\cdots<i_s=\ell}
        (-1)^{\ell-s}
        \prod_{a=0}^{s-1}b_d(p_{i_{a+1}}-p_{i_a}).
\]
This expansion follows by induction from the Hessenberg recurrence
\[
        D_m=\sum_{i=0}^{m-1}(-1)^{m-1-i}D_i\,b_d(p_m-p_i).
\]
Indeed, a chain ending at \(i\) is extended to a chain ending at \(m\) by appending the interval \([i,m]\), and the extra sign is \((-1)^{m-1-i}\).
Since \((-1)^s=(-1)^\ell(-1)^{\ell-s}\), comparing the two expressions proves the claim.
\end{proof}

\begin{proposition}\label{prop:free-product}
For odd \(d\ge5\),
\[
        \PP G_d\cong(\Z/2\Z)*\Z.
\]
\end{proposition}

\begin{proof}
Let \(F=(\Z/2\Z)*\Z\), with generators \(r\) and \(a\), where \(r^2=1\).  The natural map \(F\to\PP G_d\) sends \(r\) to the class of \(R\) and \(a\) to the class of \(A\).  We prove that its kernel is trivial.

Every nontrivial element of \(F\) is conjugate either to a nonzero pure power \(a^m\), to the involution \(r\), or to a cyclically reduced mixed word of the displayed form.  Since projective scalarity is invariant under conjugation, it is enough to treat these representatives.  The pure-power case cannot lie in the kernel: \(A^m\) is unipotent, so if it were scalar it would be the identity, but \(A\) is a nontrivial unipotent Jordan block and has infinite order.  The involution \(r\) cannot lie in the kernel either, because \(R\) has eigenvalue \(-1\) on \(\Q e_0\) and eigenvalue \(1\) on \(e_0^\perp\).

It remains to rule out a cyclically reduced mixed word.  Suppose that its matrix \(W\) were projectively scalar, say \(W=\lambda I\).  Since \(W\) has rational entries, \(\lambda\in\Q\).  Since \(W\) is a product of \(\ell\) reflections and determinant-one powers of \(A\),
\[
        \det W=(-1)^\ell.
\]
Thus
\[
        \lambda^d=(-1)^\ell.
\]
Because \(d\) is odd and \(\lambda\in\Q\), this forces
\[
        \lambda=(-1)^\ell.
\]
Therefore
\[
        Q(e_0,We_0)=2(-1)^\ell.
\]
On the other hand, Lemma~\ref{lem:reflection-det} gives
\[
        Q(e_0,We_0)=(-1)^\ell\det H_b(k_1,\ldots,k_\ell).
\]
Hence
\[
        \det H_b(k_1,\ldots,k_\ell)=2,
\]
contradicting Theorem~\ref{thm:walk-positive}, which gives \(\det H_b(k)>2\) for \(d\ge5\).  Thus no nontrivial reduced word maps to the identity in \(\PP G_d\), and the natural map is an isomorphism.
\end{proof}

\section{Density and infinite index}

For a nonisotropic vector \(u\) with \(Q(u,u)=2\), write
\[
        \rho_u(x)=x-Q(x,u)u
\]
for the corresponding reflection.

\begin{lemma}\label{lem:two-reflection-torus}
Let \(u,v\in V\) satisfy
\[
        Q(u,u)=Q(v,v)=2,
        \qquad
        |Q(u,v)|>2.
\]
Let \(H\) be the Zariski closure of a subgroup of \(\Ogrp(Q)\) containing \(\rho_u\rho_v\).  Then
\[
        X_{uv}:x\longmapsto Q(x,u)v-Q(x,v)u
\]
lies in \(\Lie(H)\).
\end{lemma}

\begin{proof}
Let \(q=Q(u,v)\) and \(P=\operatorname{span}\{u,v\}\).  The Gram determinant of \(Q|_P\) is \(4-q^2<0\), so \(P\) is a nondegenerate plane.  The element \(h=\rho_u\rho_v\) fixes \(P^\perp\) pointwise and restricts to an element of \(\SO(P,Q|_P)\).  In the basis \(u,v\), this restriction has determinant \(1\) and trace \(q^2-2\), hence it has real eigenvalues \(\mu,\mu^{-1}\) with \(\mu\ne\pm1\).  In particular, \(\mu\) is not a root of unity.

The group \(\SO(P,Q|_P)\) is a one-dimensional algebraic torus (split over \(\R\)).  A proper algebraic subgroup of a one-dimensional torus is finite, so the powers of an element with non-root-of-unity eigenvalue are Zariski dense in that torus.  Hence the identity component of the Zariski closure of \(\langle h\rangle\) has Lie algebra \(\mathfrak{so}(P,Q|_P)\), extended by zero on \(P^\perp\).  This one-dimensional Lie algebra is spanned by
\[
        x\longmapsto Q(x,u)v-Q(x,v)u.
\]
Since \(\langle h\rangle\subseteq H\), this line lies in \(\Lie(H)\).
\end{proof}

\begin{proposition}\label{prop:zariski}
The group \(G_d=\langle A,R\rangle\) is Zariski dense in \(\Ogrp(Q)\).
\end{proposition}

\begin{proof}
All Zariski closures and Lie algebras in this proof are taken after extension of scalars to \(\C\).  Since \(\Ogrp(Q)\) is defined over \(\Q\), equality of the complexified Zariski closure with \(\Ogrp(Q)_\C\) is exactly the desired rational Zariski-density statement.  For \(0\le i\le d-1\), put
\[
        v_i=A^ie_0.
\]
The group \(G_d\) contains the reflections \(\rho_{v_i}\), because
\[
        A^iRA^{-i}=\rho_{v_i}.
\]
The vectors \(v_0,\ldots,v_{d-1}\) form a basis.  Moreover,
\[
        Q(v_i,v_i)=2,
\]
and, for \(i\ne j\),
\[
        Q(v_i,v_j)=b_d(j-i)>2.
\]
Here the strict inequality follows from
\[
        \frac{1+z}{1-z}=1+2z+2z^2+\cdots.
\]
For \(m\ge1\), the coefficient of \(z^m\) in its \(d\)-th power is at least \(2d\), by choosing one of the \(d\) factors to contribute \(2z^m\) and all other factors to contribute \(1\).  Hence it is strictly larger than \(2\).

Let \(H\) be the Zariski closure of \(G_d\).  Lemma~\ref{lem:two-reflection-torus} implies that \(\Lie(H)\) contains
\[
        X_{ij}(x)=Q(x,v_i)v_j-Q(x,v_j)v_i
        \qquad(i<j).
\]
Under the standard identification
\[
        \bigwedge^2 V\cong\mathfrak{so}(Q),
        \qquad
        u\wedge v\mapsto\bigl(x\mapsto Q(x,u)v-Q(x,v)u\bigr),
\]
the \(X_{ij}\)'s form a basis of \(\mathfrak{so}(Q)\), because the \(v_i\)'s form a basis of \(V\).  Thus the identity component of \(H\) is \(\SO(Q)\).  Since \(G_d\) contains the reflection \(R\), its Zariski closure is all of \(\Ogrp(Q)\).
\end{proof}

In the basis \(e_0,\ldots,e_{d-1}\), write
\[
        \Ogrp(Q,\Z)=\{g\in\GL_d(\Z):Q(gx,gy)=Q(x,y)\text{ for all }x,y\}.
\]
Since \(b_d(n)\in\Z\), the Gram matrix of \(Q\) is integral.  The matrix of \(A\) is integral, and
\[
        R(e_j)=e_j-Q(e_j,e_0)e_0=e_j-b_d(j)e_0
\]
is integral for every \(j\).  By Proposition~\ref{prop:orthogonal-normal-form}, both \(A\) and \(R\) preserve \(Q\).  Hence
\[
        G_d=\langle A,R\rangle\leq\Ogrp(Q,\Z).
\]

\begin{lemma}\label{lem:ambient-Z2}
Every finite-index subgroup of \(\Ogrp(Q,\Z)\) contains a subgroup isomorphic to \(\Z^2\).
\end{lemma}

\begin{proof}
Let
\[
        X=\log A
        =
        \sum_{r=1}^{d-1}(-1)^{r+1}\frac{(A-I)^r}{r}.
\]
This is a rational nilpotent matrix, and \(A=\exp X\).  For every integer \(n\),
\[
        A^n=\exp(nX)
\]
preserves \(Q\).  The entries of \(\exp(tX)^tQ\exp(tX)-Q\) are polynomials in \(t\), because \(X\) is nilpotent; they vanish for all integers \(t=n\), hence vanish identically.  Differentiating at \(t=0\) gives
\[
        X\in\mathfrak{so}(Q)(\Q).
\]
Every odd power of a \(Q\)-skew operator is again \(Q\)-skew: indeed, \(Q(X^m x,y)=(-1)^m Q(x,X^m y)\).  Hence \(X^3\in\mathfrak{so}(Q)(\Q)\).  Choose a positive integer \(L\) such that
\[
        E=\exp(LX^3)
\]
has integral entries; for instance, choose \(L\) sufficiently divisible so that every matrix \(L^mX^{3m}/m!\) occurring in the finite exponential sum has integral entries.  Since \(LX^3\in\mathfrak{so}(Q)\), this exponential preserves \(Q\).  Also \(E\) is unipotent, so \(\det E=1\).  Hence the integral matrix \(E\) lies in \(\GL_d(\Z)\), and therefore
\[
        E\in\Ogrp(Q,\Z).
\]

The elements \(A=\exp(X)\) and \(E=\exp(LX^3)\) commute.  They generate a copy of \(\Z^2\).  Indeed, if
\[
        A^mE^n=I,
\]
then, since the exponents commute,
\[
        \exp(mX+nLX^3)=I.
\]
The exponent \(mX+nLX^3\) is nilpotent.  Applying the finite polynomial logarithm to \(\exp(mX+nLX^3)=I\) gives
\[
        mX+nLX^3=0.
\]
Because \(A\) is one Jordan block of size \(d\ge5\), the logarithm \(X\) is a regular nilpotent.  Equivalently, if \(N_A=A-I\), then \(X=N_A\) plus higher powers of \(N_A\), while \(X^3=N_A^3\) plus higher powers of \(N_A\).  Since \(N_A\) and \(N_A^3\) are nonzero and have different lowest powers in the \(N_A\)-adic filtration, \(X\) and \(X^3\) are linearly independent.  Hence \(m=n=0\).

Now let \(H\le\Ogrp(Q,\Z)\) have finite index.  For any \(g\in\Ogrp(Q,\Z)\), the cyclic group \(\langle g\rangle\) meets \(H\) in finite index, so some positive power of \(g\) lies in \(H\).  Applying this to \(A\) and \(E\), choose \(m,n>0\) with
\[
        A^m\in H,
        \qquad
        E^n\in H.
\]
The subgroup \(\langle A^m,E^n\rangle\) is again isomorphic to \(\Z^2\).  Thus every finite-index subgroup of \(\Ogrp(Q,\Z)\) contains \(\Z^2\).
\end{proof}

\begin{lemma}\label{lem:free-product-no-Z2}
The group \((\Z/2\Z)*\Z\) contains no subgroup isomorphic to \(\Z^2\).
\end{lemma}

\begin{proof}
The free product acts on its Bass--Serre tree with vertex stabilizers conjugate to \(\Z/2\Z\) and \(\Z\), and with trivial edge stabilizers.  An abelian subgroup acting on a tree either fixes a vertex or preserves a line.  In the first case it lies in a conjugate of a vertex stabilizer; in the second case it maps to the infinite cyclic group of translations of the preserved line, with finite kernel.  In either case it contains no \(\Z^2\).
\end{proof}

\begin{proposition}\label{prop:infinite-index}
The group \(G_d\) has infinite index in \(\Ogrp(Q,\Z)\).
\end{proposition}

\begin{proof}
By Proposition~\ref{prop:free-product},
\[
        \PP G_d\cong(\Z/2\Z)*\Z,
\]
so \(\PP G_d\) contains no \(\Z^2\).  The kernel of
\[
        G_d\longrightarrow \PP G_d
\]
is finite, since a scalar orthogonal matrix is \(\pm I\).  Therefore \(G_d\) itself contains no \(\Z^2\): otherwise its image in \(\PP G_d\) would still contain a copy of \(\Z^2\), because \(\Z^2\) is torsion-free.

If \(G_d\) had finite index in \(\Ogrp(Q,\Z)\), then Lemma~\ref{lem:ambient-Z2} would force \(G_d\) to contain a copy of \(\Z^2\), contradiction.
\end{proof}

\section{Companion coordinates}

Let \(B=RA\).  Since \(R=BA^{-1}\), we have
\[
        G_d=\langle A,R\rangle=\langle A,B\rangle.
\]
By Proposition~\ref{prop:normal-form}, there is \(P\in\GL_d(\Q)\) such that
\[
        P A P^{-1}=\Comp((x-1)^d),
        \qquad
        P B P^{-1}=\Comp((x+1)^d).
\]
Consequently
\[
        \Gamma_d^{\rm orig}:=
        \left\langle
        \Comp((x-1)^d),\Comp((x+1)^d)
        \right\rangle
        =P G_d P^{-1}.
\]
Put
\[
        Q_d^{\rm orig}=P^{-t}QP^{-1}.
\]
After multiplying by a nonzero integer, assume that \(Q_d^{\rm orig}\) has integral entries.  Multiplying a bilinear form by a nonzero scalar does not change its orthogonal group.  Then \(\Gamma_d^{\rm orig}\) preserves \(Q_d^{\rm orig}\).

We record the commensurability step explicitly.  Let
\[
        \Lambda_0=\Z^d,
        \qquad
        \Lambda_1=P^{-1}\Z^d.
\]
Then \(\Lambda_0\) and \(\Lambda_1\) are commensurable lattices in the normal-form rational vector space.  Their stabilizers
\[
        \Ogrp(Q,\Lambda_i)=\{g\in\Ogrp(Q)(\Q):g\Lambda_i=\Lambda_i\}
        \qquad(i=0,1)
\]
are commensurable arithmetic subgroups.  In particular,
\[
        \Ogrp(Q,\Lambda_1)=P^{-1}\Ogrp(Q_d^{\rm orig},\Z)P.
\]
Indeed, choose \(m\ge1\) such that
\[
        m\Lambda_0\subseteq\Lambda_1\subseteq m^{-1}\Lambda_0.
\]
The group \(\Ogrp(Q,\Lambda_0)\) acts on the finite set of lattices \(L\) with
\[
        m\Lambda_0\subseteq L\subseteq m^{-1}\Lambda_0,
\]
and the stabilizer of \(\Lambda_1\) for this action is \(\Ogrp(Q,\Lambda_0)\cap\Ogrp(Q,\Lambda_1)\).  Hence the intersection has finite index in \(\Ogrp(Q,\Lambda_0)\).  The same argument with the two lattices interchanged gives finite index in \(\Ogrp(Q,\Lambda_1)\).

The group \(G_d\) lies in both \(\Ogrp(Q,\Lambda_0)\) and \(\Ogrp(Q,\Lambda_1)\).  Indeed, the inclusion in \(\Ogrp(Q,\Lambda_0)\) was proved above, while the inclusion in \(\Ogrp(Q,\Lambda_1)\) follows from \(P G_dP^{-1}=\Gamma_d^{\rm orig}\leq\GL_d(\Z)\).  Put
\[
        K=\Ogrp(Q,\Lambda_0)\cap\Ogrp(Q,\Lambda_1).
\]
If \(\Gamma_d^{\rm orig}\) had finite index in \(\Ogrp(Q_d^{\rm orig},\Z)\), then conjugating by \(P^{-1}\) would make \(G_d\) finite index in \(\Ogrp(Q,\Lambda_1)\).  Since
\[
        G_d\subseteq K\subseteq\Ogrp(Q,\Lambda_1),
\]
we would have
\[
        [K:G_d]\leq[\Ogrp(Q,\Lambda_1):G_d]<\infty.
\]
Because \(K\) has finite index in \(\Ogrp(Q,\Lambda_0)\), it would follow that
\[
        [\Ogrp(Q,\Lambda_0):G_d]
        =
        [\Ogrp(Q,\Lambda_0):K][K:G_d]<\infty,
\]
contradicting Proposition~\ref{prop:infinite-index}.  Thus infinite index transfers to companion coordinates.  Zariski density and the projective free-product identification transfer immediately under rational conjugacy.

\begin{theorem}\label{thm:main}
For every odd \(d\ge5\), the odd all-half hypergeometric group
\[
        \Gamma_d^{\rm orig}
        =
        \left\langle
        \Comp((x-1)^d),\Comp((x+1)^d)
        \right\rangle
\]
is thin in its ambient integral orthogonal group.  More precisely, for the integral rational form \(Q_d^{\rm orig}\) constructed above,
\[
        \overline{\Gamma_d^{\rm orig}}^{\,Zar}
        =
        \Ogrp(Q_d^{\rm orig})
\]
and
\[
        [\Ogrp(Q_d^{\rm orig},\Z):\Gamma_d^{\rm orig}]=\infty.
\]
Moreover,
\[
        \PP\Gamma_d^{\rm orig}\cong(\Z/2\Z)*\Z.
\]
\end{theorem}

\begin{proof}
This is exactly the combination of Proposition~\ref{prop:free-product}, Proposition~\ref{prop:zariski}, Proposition~\ref{prop:infinite-index}, and the commensurability transfer above.
\end{proof}

\begin{remark}
This is an odd-dimensional orthogonal tower:
\[
        d=5,7,9,\ldots.
\]
It is the orthogonal analogue of the even all-half symplectic tower.  The decisive simplification in odd dimension is that the rank-one local monodromy is a reflection, so the reduced-word determinant has a single fixed subdiagonal value rather than an entire transvection-exponent cube.
\end{remark}

\end{document}
